\chapter{Setting up the SHiP software}
\label{ch:setup}

\chaptergoals{1 of 11 in Part~II.}
{Part~I, in particular \cref{ch:pixi}; a Linux x86\_64 machine with enough disk space.}
{install the SHiP software chain in one pixi environment, check that it works, and have the tutorial input files at hand.}

\begin{organisernote}
This chapter is a skeleton for the setup introduction given at the start of
the tutorial. Replace the \tbd{} items once the tutorial environment is
frozen. Whatever the final recipe, participants must leave this chapter with
\emph{one} environment in which \cmd{phlex}, \cmd{jsonnet}, the \pkg{aegir},
\pkg{Shannon} and \pkg{sea-cucumber} plugins and executables, \cmd{gmex}
and PyROOT all work, and with the tutorial input files on disk.
\end{organisernote}

\section{What you need}

\begin{itemize}
  \item A Linux x86\_64 machine: the SHiP packages are built for
        \texttt{linux-64} only. \tbd{lxplus / institute machines / provided VMs;
        macOS users?}
  \item About \tbd{N}\,GB of free disk space for the environment and the
        tutorial files.
  \item \href{https://pixi.sh}{pixi} (\cref{ch:pixi}), the package manager
        used by the new-stack repositories. It installs everything else,
        compilers included, without root access.
  \item A web browser, for the \pkg{sea-cucumber} web display (\cref{ch:display}).
\end{itemize}

\section{Where the software comes from}

Released components of the new stack are distributed as conda packages on the
\href{https://prefix.dev/channels/ship}{\texttt{prefix.dev/ship}} channel;
other dependencies are resolved from the SHiP channel and
\texttt{conda-forge}. The recipes are in \repo{ship-conda-recipes}. Each
repository carries a \file{pixi.toml} that pins its dependencies and a
\file{pixi.lock} that freezes the exact versions, so the same committed lock file can reproduce a resolved environment.
Use the tutor-provided revision and lock file; cloning moving branches does
not freeze the software. \cmd{pixi install --locked} prevents silent re-resolution.

\begin{note}
As in \cref{ch:pixi}, two commands do almost everything:
\begin{description}[style=nextline]
  \item[\cmd{pixi run <task or command>}] runs one command inside the
    environment, e.g.\ \cmd{pixi run build} or \cmd{pixi run phlex -c \ldots}.
  \item[\cmd{pixi shell}] opens a shell with the environment active until you
    \cmd{exit}. This is the convenient mode for this tutorial.
\end{description}
\end{note}

\section{The tutorial environment}

\tbd{Final recipe. Either a single tutorial \file{pixi.toml} pulling in all
packages at the versions in \cref{tab:versions}, or the per-repository
checkouts below.}

If you work with the individual repositories, check them out side by side;
the SHiP developer guide calls this the \emph{sibling layout}, and some
configuration files refer to each other by relative path:

\begin{shell}
cd "$SHIP_TUTORIAL"
for r in Geometry aegir Shannon sea-cucumber field_service; do
  git clone https://github.com/ShipSoft/$r.git
done
(cd Geometry && pixi run geometry)     # builds the detector description
(cd aegir    && pixi run build)        # builds the simulation plugins
(cd Shannon  && pixi run build)        # builds the digitisation plugins
# Before building sea-cucumber, pin shipdatamodel to >=0.5,<0.6
# in its pixi.toml and freeze the resolved environment with the tutors.
(cd sea-cucumber && pixi run build)    # after that compatibility check
\end{shell}

Inside \pkg{aegir} and \pkg{Shannon}, the pixi activation script
(\file{activate.sh}) sets the environment variables the framework needs:

\begin{center}
\begin{tabularx}{\linewidth}{>{\ttfamily}lX}
\toprule
\normalfont\bfseries Variable & \bfseries Meaning \\
\midrule
PHLEX\_PLUGIN\_PATH & where \cmd{phlex} looks for plugins: the package's
  \file{build/} directory, then the environment's \file{lib/} \\
SHIPGEOMETRY\_ROOT & a bare geometry file name such as
  \file{ship\_geometry.db} is looked up in the current directory, then in
  \file{\$SHIPGEOMETRY\_ROOT/share/geometry/} \\
SHIPFIELD\_ROOT & a bare field-map file name is looked up in
  \file{\$SHIPFIELD\_ROOT/share/field/} \\
\bottomrule
\end{tabularx}
\end{center}

\section{Versions used in this booklet}

\begin{table}[h]
\centering
\begin{tabular}{lll}
\toprule
Package & conda package & Version \\
\midrule
\pkg{Geometry}         & \texttt{shipgeometry}         & \tbd{x.y} \\
\pkg{geometry service} & \texttt{shipgeometryservice}  & \tbd{x.y} \\
\pkg{field service}    & \texttt{field-service}        & \tbd{x.y} \\
\pkg{data model}       & \texttt{shipdatamodel}        & 0.5.x \\
\pkg{ship-random}      & \texttt{ship-random}          & \tbd{x.y} \\
\pkg{aegir}            & \texttt{aegir}                & \tbd{x.y} \\
\pkg{Shannon}          & \texttt{shannon}              & \tbd{x.y} \\
\pkg{sea-cucumber}     & (built from source)           & \tbd{commit} \\
Phlex                  & \texttt{phlex}                & $\geq 0.4.1$ \\
\bottomrule
\end{tabular}
\caption{Target version constraints and values still to freeze. This is not a tested environment manifest.}
\label{tab:versions}
\end{table}

\begin{pitfall}
The data model is at the centre of a version chain: \pkg{aegir} writes it,
\pkg{Shannon} and \pkg{sea-cucumber} read it, and your analysis reads it
too. \texttt{shipdatamodel} 0.6 renamed camelCase data-model members to snake\_case, for example
\field{pdgCode} to \field{pdg\_code}. This booklet uses the 0.5 names, which is what \pkg{aegir}
and \pkg{Shannon} are built against at the time of writing.
\Cref{app:datamodel} lists both spellings. If a script fails with an
unknown member, check which data-model version your environment has before
anything else.
\end{pitfall}

\begin{organisernote}
Data model 0.6.0 (13~October~2026) is the breaking snake\_case rename.
\pkg{aegir} and \pkg{Trout} pin \texttt{shipdatamodel >=0.5,<0.6},
\pkg{Shannon} pins \texttt{>=0.3,<0.6}, and \pkg{sea-cucumber} takes
\texttt{"*"} while its source still uses the camelCase names, so a fresh
\pkg{sea-cucumber} environment may resolve 0.6 and fail to build. Freeze all
packages on one data-model version for the tutorial; if that becomes 0.6,
the plugins and analysis scripts need an API migration, and existing
digitised files need an explicit compatibility strategy. The member mapping is in
\cref{app:datamodel}.
\end{organisernote}

\begin{pitfall}
As checked on 15 October 2026, \pkg{field service}'s development manifest
requires \texttt{shipdatamodel >=0.6.0}, while \pkg{aegir} requires
\texttt{>=0.5,<0.6}. These constraints cannot be satisfied in one environment.
Use a compatible released field-service build with the simulation, or keep
field-map tools in a separate process and verify their file format against
that build. Do not add the current field-service checkout to the simulation
environment without resolving this conflict. \pkg{sea-cucumber}'s unbounded
data-model dependency also needs a compatible pin before a fresh solve.
\end{pitfall}

\section{Tutorial input files}

Two kinds of input are provided rather than generated during the tutorial:

\begin{itemize}
  \item \textbf{HNL samples}: events from \pkg{EventCalc}, already converted to
        the SHiP data model (\cref{ch:generator}). Location: \tbd{path or URL}.
  \item \textbf{Field maps}: the muon-shield and spectrometer maps in the SHiP
        field-map format (\cref{sec:fieldservice}). Location: \tbd{path or URL}.
\end{itemize}

Copy them to \file{\$SHIP\_TUTORIAL/work/}. Keep the sample metadata
(mass, lifetime, mixing, generator revision, coordinate transform, weights
and exposure) beside the files.

\begin{exercise}{Check your environment}
From the \file{aegir} directory, open a shell in the environment and check
that the tools respond:
\begin{shell}
pixi shell
which phlex jsonnet
# In the Geometry environment, check gmex separately: which gmex
echo $PHLEX_PLUGIN_PATH
python -c "import ROOT; print(ROOT.gROOT.GetVersion())"
\end{shell}
\end{exercise}

\begin{checkpoint}
Paths for Phlex and Jsonnet inside the pixi environment, a plugin path that starts with the
\file{aegir/build} directory, and a ROOT version satisfying the frozen environment (aegir currently requires
6.40 or later). Geometry viewing additionally needs \cmd{gmex} and a graphical
session; do not assume it is installed in the aegir environment.
\end{checkpoint}
